\chapter{G. S. Ghurye: Features of Caste and PYQs}

\section{Caste as Integrative and Divisive}

Ghurye held two seemingly opposite views of caste: the \textbf{traditional avatar of caste integrated} Indian society, while the \textbf{modern avatar of caste (created by the British through the caste census and quota politics) divides} it.

\begin{KeyConcept}
\begin{itemize}
\item Viewed through the \textbf{Indological perspective}, caste is part of Hindu culture and \textbf{integrates} us (like the joint family, caste structure, the triad, Vedic knowledge --- a composite, syncretic culture diffused all over India).
\item Viewed through the \textbf{westernized/rational/Enlightened} lens, caste \textbf{divides} us.
\item Giving both accounts at once is itself the \textbf{evolution of caste} --- and what worried Ghurye most was the \textbf{politicization of caste}.
\end{itemize}
\end{KeyConcept}

\begin{KeyConcept}
\begin{itemize}
\item Ghurye was the \textbf{first scholar in India to discuss the politicization of caste} --- but the \emph{phrase} \textbf{'politicization of caste'} is coined by \textbf{Rajni Kothari} (one of India's greatest political scientists).
\item Caste census + quota politics + caste associations + politicization of caste are, for Ghurye, \textbf{unhealthy for Indian society} --- and all of these came into existence only because of the British.
\end{itemize}
\end{KeyConcept}

\begin{KeyConcept}
\begin{itemize}
\item Ghurye coined the term \textbf{caste patriotism}: when the census counts you as a member of a caste, you become \textbf{conscious of your caste membership}, caste associations grow stronger, and a 'caste patriotism' develops --- what \textbf{M. N. Srinivas} later called the \textbf{vote bank} (derived from Ghurye's caste patriotism).
\item \textbf{Caste patriotism integrates the caste but divides India} --- mechanical solidarity becoming the biggest hurdle on the path to organic solidarity.
\end{itemize}
\end{KeyConcept}

\begin{ExampleBox}
\begin{itemize}
\item Caste-based regiments in the Indian Army (e.g. the \textbf{Mahar Regiment}, for which \textbf{Ambedkar} fought) reflect the solidarity around caste ties of traditional society --- and this caste-based regimentation was done by the \textbf{British}.
\item Caste was an Indo-Gangetic phenomenon diffused all over India \emph{through the caste census} --- the British treated caste as pan-Indian and imposed it even on people who had no caste.
\end{itemize}
\end{ExampleBox}

\section{The Six (Structural) Features of Caste}

\begin{KeyConcept}
\begin{itemize}
\item In \textbf{Caste and Race (1932)}, chapters 1--2, Ghurye discusses the \textbf{features of caste as it was present in the 1920s} (traditional India, before fundamental rights). These are called the \textbf{structural features of caste}.
\item The six features: (1) \textbf{segmentation of society}, (2) \textbf{hierarchy}, (3) \textbf{civil and religious disabilities and privileges}, (4) \textbf{lack of unrestricted choice of occupation}, (5) \textbf{restriction on food, drinks and social intercourse} (commensality), and (6) \textbf{endogamy}.
\end{itemize}
\end{KeyConcept}

\section{Feature 1 --- Segmentation of Society}

\begin{KeyConcept}
\begin{itemize}
\item \textbf{Segmentation of society} = stratification into segments (like the floors of a building). The tribal model speaks of \textbf{egalitarianism} (no high/low tribe), but caste has \textbf{high caste and low caste} --- a segmentation.
\item You are \textbf{born into} a caste (you do not choose it), so caste is \textbf{closed strata} based on \textbf{ascriptive features}; therefore there is \textbf{no mobility} in caste.
\item Differences between castes are not only in \textbf{rank} --- they are also reflected in \textbf{food, behaviour, dress code, culture and custom}: every caste has its own regulations and customs making it distinct.
\end{itemize}
\end{KeyConcept}

\begin{ExampleBox}
\begin{itemize}
\item \textbf{Andre Beteille's field study (1960)} --- a 10-month study of \textbf{Shripuram} village in \textbf{Tanjore (Thanjavur)} district, Tamil Nadu --- exposed the \textbf{segments within Brahmins} (a 'society' can be a small village).
\item He found four categories within Brahmins; two of them: the \textbf{Smarta Brahmins} (followers of \textbf{Shankaracharya}, worshipping Shiva and Vishnu) and the \textbf{Sri Vaishnava Brahmins} (followers of \textbf{Ramanuja}, worshipping Vishnu only).
\item The Smarta do not marry the Sri Vaishnava --- \textbf{endogamy within Brahmins} --- and the two are untouchable to each other. This shows \textbf{untouchability existed at the upper level too}, not only at the bottom.
\end{itemize}
\end{ExampleBox}

\section{Feature 2 --- Hierarchy}

\begin{KeyConcept}
\textbf{Hierarchy} = a \textbf{ranked arrangement}. Your rank in the hierarchy decides your \textbf{life chances} --- occupation, wealth, status.
\end{KeyConcept}

\section{Feature 3 --- Civil and Religious Disabilities and Privileges}

\begin{KeyConcept}
\begin{itemize}
\item \textbf{Civil and religious disabilities and privileges} --- in the hierarchy formed by segmentation, the \textbf{upper castes} get privileges and the \textbf{lower castes / casteless (untouchables)} suffer disabilities.
\item Disabilities imposed on the low caste/untouchable: they cannot drink from the common well, cannot enter the temple, cannot access education, cannot sit and dine with upper castes, and must maintain social distance.
\item These \textbf{civil and religious disabilities developed from the concept of purity and pollution}: special powers and rights are bestowed on higher castes, while many disabilities are imposed on low castes and untouchables.
\end{itemize}
\end{KeyConcept}

\section{Feature 4 --- Lack of Unrestricted Choice of Occupation}

\begin{KeyConcept}
\begin{itemize}
\item \textbf{Lack of unrestricted choice of occupation} = restrictions on the choice of occupation. Occupation is fixed by \textbf{heredity}: your caste decides your occupation, and you cannot change it --- because to change occupation means to change caste, and caste cannot be changed.
\item Occupations were fixed by heredity; individuals were traditionally not allowed to change their traditional occupation.
\item The upper castes (Brahmins) were free to opt for the study of religious books, which other caste groups could not do.
\end{itemize}
\end{KeyConcept}

\section{Feature 5 --- Restriction on Food, Drinks and Social Intercourse (Commensality)}

\begin{KeyConcept}
\begin{itemize}
\item \textbf{Restriction on food, drinks and social intercourse} --- in one word, \textbf{commensality} --- the prohibition on inter-dining and intercourse.
\item Traditional Hindu society distinguishes \textbf{kaccha food} (raw/uncooked --- e.g. fruits; any food made in \textbf{water}) from \textbf{pakka food} (cooked --- any food made in \textbf{ghee}).
\item An upper-caste person would \textbf{not accept pakka food from a low-caste person} (he might accept kaccha/raw food, whose taste the other's touch cannot spoil) --- and sometimes even water or kaccha food is not accepted.
\item \textbf{What is accepted, by whom, and from whom, reflects hierarchy and caste consciousness.} These practices are still prevalent in villages.
\end{itemize}
\end{KeyConcept}

\section{Feature 6 --- Endogamy and Marriage Restrictions}

\begin{KeyConcept}
\begin{itemize}
\item \textbf{Endogamy} = marrying within one's caste. But caste is not homogeneous: \textbf{Maheshwari} and \textbf{Agarwal} are not castes but \textbf{sub-castes} of the \textbf{Baniya} caste --- so 'caste endogamy' is really \textbf{sub-caste endogamy} (an Agarwal marrying a Maheshwari is sub-caste exogamy).
\item The unit of caste is paradoxical: \textbf{Varna was discovered by Orientalists} (their reading of the Dharma Shastras, Gita and Puranas), \textbf{Jati was discovered through the census questionnaires by British administrators}, and \textbf{caste/sub-caste as a functioning system was manufactured by social anthropologists / structural-functional field workers}.
\end{itemize}
\end{KeyConcept}

\begin{KeyConcept}
\begin{itemize}
\item \textbf{Uma Chakravarti}: the category \textbf{class is about production} (economic), while \textbf{caste is about reproduction} (social, not merely economic) --- 'class produces, caste reproduces'. Caste is reproduced only through \textbf{endogamy}.
\item \textbf{Arranged marriage} is really an arrangement \textbf{between castes} --- invented as an apparatus to reproduce castes and maintain the caste structure.
\item The \textbf{caste panchayat} (e.g. khap panchayat) was the instrument that ensured endogamous marriage.
\end{itemize}
\end{KeyConcept}

\begin{KeyConcept}
\begin{itemize}
\item \textbf{Hypergamy} (from the woman's perspective, marrying a person of same or higher caste/status --- \textbf{anuloma}) was allowed but \textbf{with restrictions}. \textbf{Hypogamy} (the husband is lower --- \textbf{pratiloma}) was discouraged.
\item According to Ghurye, marriage restrictions in society are of \textbf{three types}: \textbf{endogamy}, \textbf{exogamy}, and \textbf{hypergamy}.
\item \textbf{Exogamy} operates at the village level and the gotra level (\textbf{sapinda} --- a boy and girl with a common ancestor/'pind' cannot marry each other).
\end{itemize}
\end{KeyConcept}

\section{Ghurye on Caste Patriotism, Reservation and Politics}

\begin{CriticalPoint}
\begin{itemize}
\item \textbf{Ghurye opposed reservation} --- reservation creates caste consciousness, divides society and strengthens caste.
\item For Ghurye, the spirit of \textbf{caste patriotism can never develop patriotism cutting across caste lines}: it only reduces the nation, in our minds, to the level of caste (your caste becomes your nation). Reservation would \textbf{tear asunder national integration} --- the very vision of the constitution-makers.
\item Ghurye would therefore also \textbf{oppose the Poona Pact / separate electorate}, \textbf{EWS reservation}, and the \textbf{Fifth/Sixth Schedule} (he was an \textbf{assimilationist} --- tribals are 'backward Hindus', so do not isolate them; their gods and our gods are the same, via diffusion and assimilation --- which Srinivas would call parochialization and universalization).
\end{itemize}
\end{CriticalPoint}

\begin{ComparisonBox}
\textbf{Ghurye and Ambedkar never met}: Ghurye was a professor who limited himself to academic discourse, while Ambedkar was much more --- a practising lawyer, activist, and leader. (In the 1920s Ambedkar was also active in academics, but their paths did not cross.)
\end{ComparisonBox}

\section{Caste Outside India}

\begin{KeyConcept}
\begin{itemize}
\item Caste is \textbf{not unique to India}. It is found in Sri Lanka (traditional Ceylon), Nepal, Bhutan, Pakistan and Afghanistan. \textbf{Max Weber} also discussed caste outside India.
\item \textbf{Isabel Wilkerson} recently published a popular book on \textbf{caste in America} (after the Cisco case), where she \textbf{equated race with caste}.
\item Caste is a \textbf{political, academic and Orientalist category} developed by the British in their attempt to understand Indian culture --- classifying Indians (like we classify animals as endangered/extinct).
\item Ghurye's point: \textbf{caste is a social group having these features} --- so if these features are found outside India, those groups are also caste.
\end{itemize}
\end{KeyConcept}

\section{PYQ Content: Ambedkar, Caste and Democracy}

\begin{KeyConcept}
\begin{itemize}
\item \textbf{Ambedkar on caste}: the two most important features of the caste system are the \textbf{division of labourers} and \textbf{graded inequality} (key words), followed by \textbf{endogamy}.
\item \textbf{``Caste has been unhealthy for the growth of the sociology of India''}: \textbf{Beteille} said caste has been the \textbf{obsession of Indian anthropologists more than of the people of India} (over-obsession with any one thing is detrimental); \textbf{A. R. Desai} called caste the \textbf{``steel frame of Hinduism''} (Hinduism equated with caste). Indian sociology was born of \textbf{Indology} --- the study of Hindu culture, i.e. the study of caste.
\end{itemize}
\end{KeyConcept}

\begin{KeyConcept}
\textbf{``Has caste hindered democracy?''} --- discuss both dimensions. \textbf{Yes}: Ambedkar held \textbf{Hinduism and democracy are opposite} (democracy = equality, caste = inequality); \textbf{Ghurye} held democracy needs collective consciousness around national identity, while caste associations and caste patriotism hinder national integration. \textbf{No}: democracy itself gave castes the chance to consolidate (caste associations/pressure groups were born of democracy; elites circulate on caste lines).
\end{KeyConcept}

\begin{KeyConcept}
\textbf{Positional vs structural change}: \textbf{Sanskritization is positional change} --- adopting upper-caste customs (vegetarianism, the sacred thread, Vedic knowledge) moves you up within the caste structure but does \textbf{not erode the structure} (the status quo is maintained --- a moving equilibrium). Structural change would erode the structure itself.
\end{KeyConcept}

\begin{ExampleBox}
\textbf{The abolition of sati}: \textbf{William Bentinck} brought the legislation, supported by \textbf{Raja Ram Mohan Roy}. It was a \textbf{heterogenetic change} --- the force of modernization/westernization promoted \textbf{humanitarianism} (a feature of Westernization), which eroded the traditional institutions against the fundamental rights of Dalits, women and weaker sections. Sati was practised in the \textbf{upper castes}; through the interactionist lens, a woman's self-immolation was a \textbf{symbol of the king's pride} --- a \textbf{patriarchal bargain} (wives competing to maintain the structure of private patriarchy).
\end{ExampleBox}

\begin{KeyConcept}
\begin{itemize}
\item \textbf{A. R. Desai's Marxist perspective} (PYQ): Indian nationalism analysed through the \textbf{mode of production} --- feudal before colonialism, \textbf{colonial/capitalist} under the British, and never socialist after independence; his books are \textbf{``Recent Trends of Indian Nationalism''} and \textbf{``Path of Development''}. Critique: it is an \textbf{ideal type mistaken for actual type} (Weber), \textbf{armchair philosophy} (not empirical), and \textbf{too myopic/economistic} (everything is 'bourgeois' --- constitution, planning, cooperatives).
\item \textbf{Ranjit Guha} (PYQ): the first person to introduce \textbf{subaltern studies} --- the most unique, indigenous study, reading history from below (through the eyes of tribals and lower castes, not the British or upper castes).
\item \textbf{Critique of the structural-functional approach} (PYQ): the general critique is the \textbf{Marxist perspective}, and the second is its \textbf{over-emphasis on integration and neglect of conflict}.
\end{itemize}
\end{KeyConcept}

\begin{QuickRevision}
\begin{itemize}
\item Ambedkar: division of labourers + graded inequality + endogamy.
\item Caste unhealthy for sociology: Beteille (obsession) + A. R. Desai (steel frame of Hinduism); sociology born of Indology.
\item Caste \& democracy: Ambedkar/Ghurye (hinders) vs democracy consolidating caste.
\item Positional (Sanskritization) vs structural change.
\item Sati: Bentinck, heterogenetic, humanitarianism, private patriarchy, patriarchal bargain.
\item A. R. Desai (mode of production, ideal vs actual type, armchair, economistic); Ranjit Guha (subaltern); structural-functional critique (Marxist, neglect of conflict).
\end{itemize}
\end{QuickRevision}

\section{Conclusion and the PYQ Approach}

\begin{KeyConcept}
\begin{itemize}
\item In conclusion, \textbf{Ghurye's account of caste laid the foundation upon which the superstructure of Indian sociology got erected} --- because Indian sociology has more or less been the sociology of caste.
\item There is a natural chain of critique: \textbf{Ghurye $\rightarrow$ M. N. Srinivas (a critique of Ghurye) $\rightarrow$ Louis Dumont (a critique of Srinivas) $\rightarrow$ Andre Beteille (a critique of Dumont).}
\end{itemize}
\end{KeyConcept}

\begin{QuickRevision}
\begin{itemize}
\item Ghurye: traditional caste integrates, modern (census/quota) caste divides; politicization of caste (phrase by Rajni Kothari).
\item Caste census $\rightarrow$ caste consciousness $\rightarrow$ caste associations $\rightarrow$ caste patriotism (Ghurye) $\rightarrow$ vote bank (Srinivas); caste patriotism integrates caste, divides India.
\item Six structural features (Caste and Race, ch. 1--2, 1920s): segmentation, hierarchy, civil/religious disabilities \& privileges, lack of unrestricted choice of occupation, restriction on food/drink/intercourse (commensality, kaccha vs pakka), endogamy.
\item Segmentation example: Beteille's Shripuram/Tanjore (1960) --- Smarta vs Sri Vaishnava Brahmins (endogamy + untouchability at upper level).
\item Endogamy: sub-caste endogamy (Maheshwari/Agarwal = Baniya sub-castes); Varna (Orientalists) vs Jati (census) vs caste (anthropologists); Uma Chakravarti (caste reproduces); arranged marriage; caste panchayat; anuloma/pratiloma; three restrictions (endogamy/exogamy/hypergamy); sapinda/gotra exogamy.
\item Ghurye opposed reservation, Poona Pact, EWS, Fifth/Sixth Schedule (assimilationist).
\item Caste outside India: Sri Lanka, Nepal, Bhutan, Pakistan, Afghanistan; Isabel Wilkerson's 'Caste'; Max Weber.
\item Conclusion: Ghurye's account = foundation of Indian sociology; chain Ghurye$\rightarrow$Srinivas$\rightarrow$Dumont$\rightarrow$Beteille.
\end{itemize}
\end{QuickRevision}

